\documentclass[12pt,a4paper]{article}
\usepackage[utf8]{inputenc}
\usepackage[T1]{fontenc}
\usepackage{geometry}
\usepackage{booktabs}
\usepackage{hyperref}
\usepackage{listings}
\usepackage{xcolor}
\usepackage{amsmath}
\usepackage{graphicx}
\usepackage{enumitem}

\geometry{margin=2.5cm}

\hypersetup{
    colorlinks=true,
    linkcolor=blue,
    urlcolor=blue,
    citecolor=blue
}

\lstset{
    basicstyle=\ttfamily\small,
    breaklines=true,
    frame=single,
    backgroundcolor=\color{gray!10}
}

\title{
    \textbf{CENG 467 --- Natural Language Understanding and Generation} \\
    \large Progress Report: Alignment via Direct Preference Optimization (DPO) \\
    \normalsize İzmir Institute of Technology, Spring 2026
}

\author{
    Enes Ergün Hoşgör (300201053) \\
    Samet Buldanlıoğlu (300201085) \\
    Uğur Mert Çavuşoğlu (300201087)
}

\date{May 14, 2026}

\begin{document}

\maketitle

% ─────────────────────────────────────────────────────────────
\section{Project Title, Group Members, and Problem Definition}

\textbf{Project Title:} Alignment via Direct Preference Optimization (DPO)

\subsection{Problem Definition}

Large language models (LLMs) trained on internet-scale corpora tend to generate outputs that are
unhelpful, harmful, or misaligned with human values. The task of \textit{alignment} aims to steer
model behavior so that generated responses are helpful, harmless, and honest (HHH).

Classical alignment pipelines rely on Reinforcement Learning from Human Feedback (RLHF) with
Proximal Policy Optimization (PPO), which requires training a separate reward model and is known
to be unstable. \textbf{Direct Preference Optimization (DPO)} \cite{rafailov2023} reformulates the
RLHF objective as a supervised classification problem over preference pairs, eliminating the need
for a reward model entirely.

In this project, we implement and evaluate DPO on the \texttt{Anthropic/hh-rlhf} dataset---a
large-scale human preference dataset containing pairs of helpful/harmless assistant responses.
The input is a multi-turn conversation prompt; the output is a natural language response. The core
challenge is teaching the model to prefer safe, helpful responses over harmful or evasive ones,
without catastrophic forgetting of language modeling ability.

% ─────────────────────────────────────────────────────────────
\section{Dataset / Benchmark Status}

\begin{itemize}[leftmargin=*]
    \item \textbf{Dataset:} Anthropic HH-RLHF \cite{bai2022}
    \item \textbf{Source:} \url{https://huggingface.co/datasets/Anthropic/hh-rlhf}
    \item \textbf{Size:} 160,800 training samples, 8,505 test samples
    \item \textbf{Format:} Each sample contains a \texttt{chosen} and a \texttt{rejected} field,
          both containing full conversation histories in the format:\\
          \texttt{\textbackslash n\textbackslash nHuman: ...\textbackslash n\textbackslash nAssistant: ...}
    \item \textbf{Split:} Train / Test (no official validation split; we use test for evaluation)
\end{itemize}

\subsection{Preprocessing}

The raw dataset contains full conversation strings. We extract the last assistant turn to separate
\texttt{prompt}, \texttt{chosen} response, and \texttt{rejected} response:

\begin{lstlisting}[language=Python]
def extract_prompt_and_response(sample):
    prompt = sample["chosen"].rsplit("\n\nAssistant:", 1)[0] + "\n\nAssistant:"
    chosen = sample["chosen"].rsplit("\n\nAssistant:", 1)[1].strip()
    rejected = sample["rejected"].rsplit("\n\nAssistant:", 1)[1].strip()
    return {"prompt": prompt, "chosen": chosen, "rejected": rejected}
\end{lstlisting}

After filtering malformed samples (empty responses, identical chosen/rejected pairs), the final
dataset contains 159,876 training and 8,505 test samples. The dataset has been fully downloaded,
preprocessed, and verified locally and on Colab.

% ─────────────────────────────────────────────────────────────
\section{Literature Review Progress}

\subsection{Key Papers Reviewed}

\textbf{Rafailov et al. (2023) --- DPO \cite{rafailov2023}:}
The foundational paper for this project. DPO shows that the optimal RLHF policy can be expressed
in closed form, allowing alignment to be solved as a binary cross-entropy problem over preference
pairs. The DPO loss is:
\[
\mathcal{L}_\text{DPO}(\pi_\theta) = -\mathbb{E}_{(x,y_w,y_l)} \left[
    \log \sigma\!\left(
        \beta \log \frac{\pi_\theta(y_w|x)}{\pi_\text{ref}(y_w|x)}
        - \beta \log \frac{\pi_\theta(y_l|x)}{\pi_\text{ref}(y_l|x)}
    \right)
\right]
\]
where $y_w$ is the chosen response, $y_l$ the rejected response, $\pi_\text{ref}$ the frozen
reference model, and $\beta$ controls deviation from the reference policy.

\textbf{Bai et al. (2022) --- HH-RLHF \cite{bai2022}:}
Introduces the dataset we use and the HHH (helpful, harmless, honest) framework. Establishes the
preference data collection methodology used by Anthropic.

\textbf{Ouyang et al. (2022) --- InstructGPT \cite{ouyang2022}:}
The seminal RLHF-PPO alignment paper. Shows significant improvement in helpfulness and safety over
base GPT-3. We use this as the baseline for understanding the problem DPO aims to simplify.

\textbf{Hu et al. (2021) --- LoRA \cite{hu2021}:}
Low-Rank Adaptation enables fine-tuning of large models by injecting trainable rank decomposition
matrices into frozen weights. We use LoRA (rank=16) throughout to fit TinyLlama on a single GPU.

\subsection{Technical Direction}

Based on the literature, we follow the standard DPO pipeline: (1) SFT on chosen responses,
(2) DPO fine-tuning using the SFT model as reference, (3) evaluation with automatic metrics
and qualitative analysis.

% ─────────────────────────────────────────────────────────────
\section{Baseline Models}

We implement two baselines for comparison with the DPO model:

\begin{itemize}[leftmargin=*]
    \item \textbf{Baseline 1 --- Zero-Shot:} The pretrained TinyLlama-1.1B model with no
          fine-tuning. Represents the unaligned lower bound. Responds to harmful prompts without
          refusal and exhibits repetitive generation patterns.

    \item \textbf{Baseline 2 --- SFT:} TinyLlama-1.1B fine-tuned with LoRA (r=16) on the
          \texttt{chosen} responses only using supervised next-token prediction. No preference
          optimization is applied. Trained for 2 epochs on 5,000 samples using
          \texttt{TRL SFTTrainer}.
\end{itemize}

Both baselines have been fully implemented and evaluated. Initial results are shown in Section~5.

% ─────────────────────────────────────────────────────────────
\section{Initial Experimental Results}

All three models were evaluated on 100 test samples using:
\begin{itemize}[leftmargin=*]
    \item \textbf{Perplexity:} Computed as $\exp(\text{mean NLL loss})$ on the chosen test responses.
          Lower is better.
    \item \textbf{BERTScore F1:} Semantic similarity between generated responses and chosen
          references using RoBERTa-Large. Higher is better.
\end{itemize}

\begin{table}[h]
\centering
\caption{Evaluation Results (100 test samples)}
\begin{tabular}{lccc}
\toprule
\textbf{Model} & \textbf{Perplexity} $\downarrow$ & \textbf{BERTScore F1} $\uparrow$ & \textbf{Status} \\
\midrule
Zero-Shot (TinyLlama-1.1B)   & 17.28 & 0.8303 & Complete \\
SFT (5k samples, 2 epochs)   & \textbf{10.08} & 0.8234 & Complete \\
DPO ($\beta$=0.1, 3k, 1 ep.) & 17.67 & \textbf{0.8272} & Complete \\
\bottomrule
\end{tabular}
\end{table}

\subsection{Interpretation}

SFT achieves the lowest perplexity (10.08), significantly improving language modeling quality by
training on preferred responses. DPO shows slightly higher perplexity than zero-shot due to
limited training data (3,000 samples, 1 epoch) and the fact that DPO optimizes for preference
margins rather than pure language modeling likelihood. However, DPO achieves the highest
BERTScore (0.8272), indicating better semantic alignment with human-preferred responses.

Qualitatively, the most important difference is in behavior on harmful prompts. The zero-shot
model provides lock-picking instructions and endorses alcohol consumption, while the DPO model
shows greater tendency toward refusal. Full qualitative comparison will be provided in the final
report.

% ─────────────────────────────────────────────────────────────
\section{Planned Improvements and Technical Direction}

For the final submission, we plan to:

\begin{enumerate}[leftmargin=*]
    \item \textbf{Scale training data:} Increase from 5k/3k to 50k+ samples for both SFT and DPO.
    \item \textbf{More training epochs:} 2-3 epochs for DPO to allow better preference learning.
    \item \textbf{Win-rate evaluation:} Use a reward model to compute DPO vs.\ SFT win-rate.
    \item \textbf{Ablation study:} Sweep $\beta \in \{0.01, 0.05, 0.1, 0.2, 0.5\}$ and LoRA
          rank $\in \{8, 16, 32\}$.
\end{enumerate}

% ─────────────────────────────────────────────────────────────
\section{Ablation / Prompt Sensitivity Plan}

The primary ablation will vary the DPO temperature parameter $\beta$, which controls the
trade-off between alignment and deviation from the reference policy:

\begin{itemize}[leftmargin=*]
    \item $\beta \in \{0.01, 0.05, 0.1, 0.2, 0.5\}$ --- primary ablation
    \item Learning rate $\in \{1\text{e-}5, 5\text{e-}5, 1\text{e-}4\}$
    \item LoRA rank $\in \{8, 16, 32\}$
\end{itemize}

Results will be reported in a table showing perplexity and BERTScore for each configuration.

% ─────────────────────────────────────────────────────────────
\section{Error Analysis and Generation Quality Plan}

We will select 20--30 test samples where the DPO model fails to outperform the SFT baseline and
categorize failures into:

\begin{itemize}[leftmargin=*]
    \item \textbf{Hallucination:} Factually incorrect claims stated with confidence.
    \item \textbf{Refusal failure:} Model complies with clearly harmful requests.
    \item \textbf{Repetition:} Degenerate looping behavior (observed in SFT outputs).
    \item \textbf{Off-topic:} Response does not address the prompt.
\end{itemize}

Generation quality will also be analyzed for coherence, fluency, and relevance on a random
sample of 50 test prompts.

% ─────────────────────────────────────────────────────────────
\section{Ethical Considerations and Bias}

The HH-RLHF dataset introduces several ethical concerns:

\begin{itemize}[leftmargin=*]
    \item \textbf{Annotator bias:} Human preference labels reflect the demographics and values of
          Anthropic's crowd workers, which may not generalize across cultures or languages.
    \item \textbf{Over-refusal:} DPO may cause the model to refuse benign requests that superficially
          resemble harmful ones.
    \item \textbf{Reward hacking:} The model may learn to produce outputs that satisfy the DPO
          loss without being genuinely helpful or harmless.
    \item \textbf{Adversarial misuse:} An aligned model can still be misled via prompt injection
          or jailbreaks.
\end{itemize}

% ─────────────────────────────────────────────────────────────
\section{GitHub and Reproducibility Status}

\begin{itemize}[leftmargin=*]
    \item \textbf{Repository:} \url{https://github.com/ugurcavusoglu/ceng467-alignment-dpo}
    \item \textbf{Structure:} \texttt{data/}, \texttt{scripts/}, \texttt{notebooks/},
          \texttt{results/tables/}, \texttt{report/}
    \item \textbf{Scripts:} \texttt{data/prepare\_dataset.py}, \texttt{scripts/train\_sft.py},
          \texttt{scripts/train\_dpo.py}, \texttt{scripts/run\_eval.py}
    \item \textbf{Dependencies:} \texttt{requirements.txt} present
    \item \textbf{README:} Present with step-by-step reproduction instructions
    \item \textbf{Results:} \texttt{results/tables/eval\_results.json} committed
    \item \textbf{Commit history:} Active development with consistent commits
\end{itemize}

To reproduce results:
\begin{lstlisting}
git clone https://github.com/ugurcavusoglu/ceng467-alignment-dpo.git
cd ceng467-alignment-dpo
pip install transformers==4.44.2 trl==0.10.1 peft==0.12.0 accelerate==0.34.0 datasets
python scripts/train_sft.py --subset 5000
python scripts/train_dpo.py --subset 3000
python scripts/run_eval.py --n_samples 100
\end{lstlisting}

% ─────────────────────────────────────────────────────────────
\section{Current Challenges and Next Steps}

\subsection{Challenges}
\begin{itemize}[leftmargin=*]
    \item \textbf{GPU memory:} TinyLlama with LoRA fits on a T4/A100, but DPO requires loading
          both model and reference model simultaneously, increasing memory pressure.
    \item \textbf{Session limits:} Colab/Kaggle session timeouts require checkpointing strategy
          for longer training runs.
    \item \textbf{TRL API instability:} TRL v1.3.0 introduced breaking changes; we pinned to
          v0.10.1 for stability.
    \item \textbf{Limited training data:} Current results use 3--5k samples; scaling to 50k+
          is needed for final evaluation.
\end{itemize}

\subsection{Next Steps (Weeks 12--15)}
\begin{enumerate}[leftmargin=*]
    \item Scale SFT and DPO training to 50k+ samples
    \item Run ablation study over $\beta$ values
    \item Implement win-rate evaluation with a reward model
    \item Perform error analysis on 20--30 failure cases
    \item Write final LNCS paper (6--8 pages)
    \item Prepare live demo
\end{enumerate}

% ─────────────────────────────────────────────────────────────
\begin{thebibliography}{9}

\bibitem{rafailov2023}
Rafailov, R., Sharma, A., Mitchell, E., Ermon, S., Manning, C. D., \& Finn, C. (2023).
\textit{Direct Preference Optimization: Your Language Model is Secretly a Reward Model}.
arXiv:2305.18290.

\bibitem{bai2022}
Bai, Y., Jones, A., Ndousse, K., Askell, A., Chen, A., DasSarma, N., ... \& Kaplan, J. (2022).
\textit{Training a Helpful and Harmless Assistant with Reinforcement Learning from Human Feedback}.
arXiv:2204.05862.

\bibitem{ouyang2022}
Ouyang, L., Wu, J., Jiang, X., Almeida, D., Wainwright, C. L., Mishkin, P., ... \& Lowe, R. (2022).
\textit{Training language models to follow instructions with human feedback}.
NeurIPS 2022.

\bibitem{hu2021}
Hu, E. J., Shen, Y., Wallis, P., Allen-Zhu, Z., Li, Y., Wang, S., ... \& Chen, W. (2021).
\textit{LoRA: Low-Rank Adaptation of Large Language Models}.
arXiv:2106.09685.

\end{thebibliography}

\end{document}
